% !TEX program = xelatex
% 공통수학2 · Ⅱ. 집합과 명제 · 2. 명제 · 01 명제와 조건 · 1/2차시
% 학생용 활동지 (답안 없음)
\documentclass[11pt,a4paper]{article}
\usepackage{kotex}
\usepackage{iftex}
\ifPDFTeX\else
  \setmainhangulfont{Noto Serif CJK KR}
  \setsanshangulfont{Noto Sans CJK KR}
\fi
\usepackage[margin=2.1cm]{geometry}
\usepackage{parskip}
\usepackage{enumitem}
\usepackage{array}
\usepackage{tabularx}
\usepackage{xcolor}
\usepackage{amssymb}
\usepackage{tikz}
\usepackage[most]{tcolorbox}
\usepackage{fancyhdr}
\usepackage{titlesec}

\definecolor{goldc}{HTML}{B07C22}
\definecolor{goldstrong}{HTML}{8A5F16}
\definecolor{indigoc}{HTML}{2B3B57}
\definecolor{indigosoft}{HTML}{6E82A6}
\definecolor{linec}{HTML}{D6DACB}
\definecolor{surface2}{HTML}{E4E8DC}
\definecolor{writeline}{HTML}{C7CCBA}
\definecolor{inksoft}{HTML}{52604F}

\titleformat{\section}{\normalfont\sffamily\bfseries\large\color{indigoc}}{\thesection}{0.6em}{}
\titlespacing{\section}{0pt}{14pt}{6pt}
\newtcolorbox{goalbox}{enhanced, breakable, colback=white, colframe=goldc, boxrule=0.5pt, leftrule=3pt, arc=2pt, left=10pt, right=10pt, top=8pt, bottom=8pt}
\newtcolorbox{sheetbox}{enhanced, breakable, colback=white, colframe=linec, boxrule=0.6pt, arc=3pt, left=11pt, right=11pt, top=9pt, bottom=9pt}
\newcommand{\writespace}[1]{%
  \par\nobreak\vspace{3pt}
  \foreach \n in {1,...,#1} {%
    \noindent\textcolor{writeline}{\rule{\linewidth}{0.4pt}}\par\vspace{0.62cm}%
  }
}
\newcommand{\fillblank}[1]{\makebox[#1]{\hrulefill}}
\pagestyle{fancy}\fancyhf{}\renewcommand{\headrulewidth}{0pt}
\fancyfoot[L]{\footnotesize\color{inksoft} 공통수학2 · 2026학년도 2학기 · 지평선고등학교}
\fancyfoot[R]{\footnotesize\color{inksoft} 다음 시간 --- 01 명제와 조건(2)}

\begin{document}

\noindent
\begin{minipage}[b]{0.62\linewidth}
  {\sffamily\footnotesize\color{indigosoft} 공통수학2 · Ⅱ. 집합과 명제 · 2. 명제}\\[2pt]
  {\sffamily\bfseries\Large 01 명제와 조건 \textperiodcentered\ 26차시}
\end{minipage}\hfill
\begin{minipage}[b]{0.36\linewidth}
  \raggedleft\small\color{inksoft}
  반\ \fillblank{1.6cm} \quad 번호\ \fillblank{1.6cm}\\[4pt]
  이름\ \fillblank{3.6cm}
\end{minipage}

\vspace{4pt}
\noindent{\color{goldc}\rule{\linewidth}{2.2pt}}
\vspace{10pt}

\begin{goalbox}
\textbf{오늘의 목표}\\
명제와 조건의 뜻을 알고, 명제의 참·거짓과 부정을 판별하며 조건의 진리집합을 구할 수 있다.
\vspace{4pt}
\small\color{inksoft}
$\square$ 자신 있음 \quad $\square$ 조금 헷갈림 \quad $\square$ 처음 봄 \hfill (수업 전 나의 상태에 표시)
\end{goalbox}

\section{생각 열기 --- 참, 거짓, 판별불가}

\begin{sheetbox}
다음 세 문장을 참, 거짓, 판별할 수 없음으로 분류해 보자.

\begin{enumerate}[leftmargin=1.4em, itemsep=2pt]
  \item 백록담은 한라산에 있다.
  \item 한라산은 경기도에 있다.
  \item 한라산은 높은 산이다.
\end{enumerate}
\writespace{3}
\end{sheetbox}

\section{개념 정리 --- 명제와 조건}

\begin{sheetbox}
참, 거짓을 명확하게 판별할 수 있는 문장이나 식을 \textbf{명제}라고 한다.

\vspace{4pt}
명제 $p$에 대하여 `$p$가 아니다'를 $p$의 \textbf{부정}이라 하고, 기호 $\sim p$로 나타낸다.

\vspace{8pt}
변수를 포함하는 문장이나 식 중에서 변수의 값에 따라 참, 거짓을 판별할 수 있는 것을 \textbf{조건}이라 한다. 전체집합 $U$의 원소 중에서 조건 $p$가 참이 되게 하는 모든 원소의 집합을 조건 $p$의 \textbf{진리집합}이라고 한다.

\vspace{6pt}
\textbf{보기} --- $U=\{x\mid x$는 10 이하의 짝수$\}$, $p:x\ge6$일 때, 진리집합 $P=\ \fillblank{3.2cm}$, $\sim p:x<6$의 진리집합은 $P^C=\ \fillblank{3.2cm}$
\end{sheetbox}

\section{연습 문제}

\begin{sheetbox}
\textbf{1.} 다음에서 명제를 모두 찾고, 그 참·거짓을 판별하시오.\\
\hspace*{1.6em}⑴ $\pi$는 유리수이다. \quad ⑵ 섬진강은 아름답다. \quad ⑶ $4-1=3$ \quad ⑷ 정삼각형은 이등변삼각형이다.
\writespace{3}

\vspace{6pt}
\textbf{2.} 다음 명제의 부정을 말하고, 참·거짓을 판별하시오.\\
\hspace*{1.6em}⑴ 15의 약수의 개수는 3이다. \qquad ⑵ $2\le1+\dfrac12$
\writespace{3}

\vspace{6pt}
\textbf{3.} 다음 용어의 정의를 말하시오.\\
\hspace*{1.6em}⑴ 직사각형 \qquad ⑵ 공집합
\writespace{2}

\vspace{6pt}
\textbf{4.} 전체집합 $U$가 자연수 전체의 집합일 때, 조건 `$p:x^2>10$'에 대하여 $\sim p$와 그 진리집합을 구하시오.
\writespace{2}
\end{sheetbox}

\section{사고의 가시화 --- See · Think · Wonder}

\noindent{\footnotesize\color{inksoft} 오늘 배운 `명제와 조건'을 떠올리며 아래 세 칸을 채워 보자.}\\[6pt]
\begin{tabularx}{\linewidth}{@{}XXX@{}}
\textbf{\footnotesize\color{indigoc} See --- 무엇을 보았나?} &
\textbf{\footnotesize\color{indigoc} Think --- 무엇을 생각했나?} &
\textbf{\footnotesize\color{indigoc} Wonder --- 무엇이 궁금한가?} \\[4pt]
\end{tabularx}
\noindent
\begin{minipage}[t]{0.32\linewidth}\writespace{3}\end{minipage}\hfill
\begin{minipage}[t]{0.32\linewidth}\writespace{3}\end{minipage}\hfill
\begin{minipage}[t]{0.32\linewidth}\writespace{3}\end{minipage}

\section{마무리 성찰}

\begin{sheetbox}
\begin{minipage}[t]{0.48\linewidth}
{\footnotesize\color{inksoft} 오늘 배운 것을 한 문장으로}
\writespace{2}
\end{minipage}\hfill
\begin{minipage}[t]{0.48\linewidth}
{\footnotesize\color{inksoft} 감사 일기 / 유념 한 줄}
\writespace{2}
\end{minipage}
\end{sheetbox}

\end{document}
